% Statement and cited-input packet. Target proof is separately embargoed.
% Frozen II source SHA256: 74497bfe327fae25c9cdcb17334b3bf42a9793a9df7c0a4404d97568f2f760f0
% Target owner: wbII:lem:picard-finite-norm; tag: 03IE

\begin{lemma}[the finite-cover norm identity]
\hypertarget{tag-03IE}{}\label{wbII:lem:picard-finite-norm}


\ Retain the notation of
\cref{wbII:lem:picard-crossed-defect}; \CompendiumRef{I}{wb:prop:candidate-finite-fillings}.  For
$j=1,2$, let
\[
 N_{T_j}=I+T_j+\cdots+T_j^{m_j-1}.
\]
With the clockwise conventions of \eqref{wbII:eq:W:5.17g}, the pullback to the
$m_j$-fold $s_j$-cover gives
\[
 N_{T_j}d_j=\iota_{v_j}E
 \quad\text{in }V_{\mathbb R},
\qquad
 m_j\bar d_j=[\iota_{v_j}E]
 \quad\text{in }V_G\otimes\mathbb R.
\tag{W.28}\hypertarget{tag-03GV}{}\label{wbII:eq:W:5.17k}
\]
\end{lemma}

% Cited input wbII:eq:W:5.17g; statement or labelled formula only.
\begin{lemma}[the pair-of-pants defect class]
\hypertarget{tag-03IC}{}\label{wbII:lem:picard-crossed-defect}


\ Choose the compact pair-of-pants core $P\subset B^\circ$ and
the zero-section whiskers used in \CompendiumRef{I}{mapsI:top:prop:geometric-pi1}, and put
$J_P=f^{-1}(P)$.  Orient its boundary meridians clockwise.  Write
$\rho_0$ for the projection of the zero-section cusp meridian
$\widetilde a_0$.  The standard family uses the tautological cusp
gluing, whose affine regluing class is $\mu=0$ (equivalently
$\ell_0=\gamma(\mu)=0$); hence the flat-marked cusp lift has no vertical
translation.  The base loops and their chosen lifts therefore satisfy
\[
 \rho_1\rho_2\rho_0=1,
 \qquad
 \widehat\rho_1\widehat\rho_2\widetilde a_0=1.
\tag{W.23}\hypertarget{tag-03GQ}{}\label{wbII:eq:W:5.17g}
\]
For $L\in\operatorname{Pic}(X)$, write
$E=c_1(L|_F)\in\bigwedge^2V$.  Then $E\in(\bigwedge^2V)^G$, and there is a
crossed homomorphism
\[
 z:\pi_1(P)\longrightarrow V_{\mathbb R},\qquad
 z(gh)=z(g)+T_gz(h),
\tag{W.24}\hypertarget{tag-03GR}{}\label{wbII:eq:W:5.17h}
\]
whose class records
$c_1(L|_{J_P})-[\Omega_E]$.  Here $\Omega_E$ is the closed, translation-
invariant vertical two-form induced by $E$ in the flat marking.  If
$d_i=z(\rho_i)$ and a bar denotes the image in
$V_G\otimes\mathbb R$, then
\[
 \bar d_1+\bar d_2+\bar d_0=0.
\tag{W.25}\hypertarget{tag-03GS}{}\label{wbII:eq:W:5.17i}
\]
The three barred classes, and the number
$\langle d_0,\widehat\delta\rangle$, do not depend on the chosen crossed-
homomorphism representative.  More precisely, use the base-first convention
for the edge injection in \eqref{wbII:eq:W:5.17j}.  For any oriented loop $\rho\subset P$
and any $A_\rho$-invariant $\nu\in\Lambda$, its flat-marked sweep satisfies
\[
 \bigl\langle z(\rho),\nu\bigr\rangle
 =-\bigl\langle c_1(L|_{J_P})-[\Omega_E],
 a_*([S^1_\nu]\times[\widetilde\rho])\bigr\rangle.
\tag{W.26}\hypertarget{tag-03GU}{}\label{wbII:eq:W:5.17j'}
\]
Here $a$ is the fibrewise translation action and $\widetilde\rho$ is the
flat-marked lift of $\rho$.
\end{lemma}

% Cited input wbII:lem:picard-crossed-defect; statement or labelled formula only.
\begin{lemma}[the pair-of-pants defect class]
\hypertarget{tag-03IC}{}\label{wbII:lem:picard-crossed-defect}


\ Choose the compact pair-of-pants core $P\subset B^\circ$ and
the zero-section whiskers used in \CompendiumRef{I}{mapsI:top:prop:geometric-pi1}, and put
$J_P=f^{-1}(P)$.  Orient its boundary meridians clockwise.  Write
$\rho_0$ for the projection of the zero-section cusp meridian
$\widetilde a_0$.  The standard family uses the tautological cusp
gluing, whose affine regluing class is $\mu=0$ (equivalently
$\ell_0=\gamma(\mu)=0$); hence the flat-marked cusp lift has no vertical
translation.  The base loops and their chosen lifts therefore satisfy
\[
 \rho_1\rho_2\rho_0=1,
 \qquad
 \widehat\rho_1\widehat\rho_2\widetilde a_0=1.
\tag{W.23}\hypertarget{tag-03GQ}{}\label{wbII:eq:W:5.17g}
\]
For $L\in\operatorname{Pic}(X)$, write
$E=c_1(L|_F)\in\bigwedge^2V$.  Then $E\in(\bigwedge^2V)^G$, and there is a
crossed homomorphism
\[
 z:\pi_1(P)\longrightarrow V_{\mathbb R},\qquad
 z(gh)=z(g)+T_gz(h),
\tag{W.24}\hypertarget{tag-03GR}{}\label{wbII:eq:W:5.17h}
\]
whose class records
$c_1(L|_{J_P})-[\Omega_E]$.  Here $\Omega_E$ is the closed, translation-
invariant vertical two-form induced by $E$ in the flat marking.  If
$d_i=z(\rho_i)$ and a bar denotes the image in
$V_G\otimes\mathbb R$, then
\[
 \bar d_1+\bar d_2+\bar d_0=0.
\tag{W.25}\hypertarget{tag-03GS}{}\label{wbII:eq:W:5.17i}
\]
The three barred classes, and the number
$\langle d_0,\widehat\delta\rangle$, do not depend on the chosen crossed-
homomorphism representative.  More precisely, use the base-first convention
for the edge injection in \eqref{wbII:eq:W:5.17j}.  For any oriented loop $\rho\subset P$
and any $A_\rho$-invariant $\nu\in\Lambda$, its flat-marked sweep satisfies
\[
 \bigl\langle z(\rho),\nu\bigr\rangle
 =-\bigl\langle c_1(L|_{J_P})-[\Omega_E],
 a_*([S^1_\nu]\times[\widetilde\rho])\bigr\rangle.
\tag{W.26}\hypertarget{tag-03GU}{}\label{wbII:eq:W:5.17j'}
\]
Here $a$ is the fibrewise translation action and $\widetilde\rho$ is the
flat-marked lift of $\rho$.
\end{lemma}

